% Referee report on D. Joyner, "On the transfer of Shalika germs for G_2" (1994).
% Written by Claude (Anthropic), November 13, 2026.
% newtx is not installed in the build environment, so Times is provided by
% mathptmx, the fallback allowed by the house style guide.
\documentclass[12pt,letterpaper]{article}
\usepackage[T1]{fontenc}
\usepackage{mathptmx}
\usepackage{courier}
\usepackage[scaled=0.92]{helvet}
\usepackage[letterpaper,margin=1in]{geometry}
\usepackage{amsmath,amssymb,amsthm}
\usepackage{longtable,booktabs,array}
\usepackage{enumitem}
\usepackage{microtype}
\usepackage[hidelinks]{hyperref}
\setlength{\emergencystretch}{3em}
\setlist{itemsep=2pt,topsep=4pt}

\newcommand{\OF}{\mathcal{O}_F}
\newcommand{\ovF}{\overline{F}}
\newcommand{\cO}{\mathcal{O}}
\DeclareMathOperator{\Ind}{Ind}
\newcommand{\Lv}[1]{\textsf{L#1}}
\newcommand{\verdict}[1]{\textbf{Verdict:} \emph{#1}.}
\newcommand{\level}[1]{\textbf{Level:} #1.}
\newcommand{\loc}[1]{\textbf{Where:} #1.}
\newcommand{\fix}[1]{\textbf{Suggested repair:} #1}

\theoremstyle{definition}
\newtheorem{issue}{Issue}[section]

\title{Referee report on\\
``On the transfer of Shalika germs for $G_2$''\\
by David Joyner (manuscript, summer 1994)}
\author{Referee: Claude (Anthropic), an AI system}
\date{November 13, 2026}

\begin{document}
\maketitle

\begin{abstract}
This report was written by an AI system (Claude, made by Anthropic) at the
author's request. It is not an independent certification of the results. The
manuscript proposes to prove the matching of the $1$-, $2$- and $3$-regular
Shalika germs of the split exceptional group $G_2$ with germs on its cuspidal
endoscopic groups, and hence local transfer for $G_2$. The geometric part of the
paper is largely sound. I verified by computer that the variety $Z$ of
\S3 is normal, that the essential patches of \S4 are regular and cover $Z$ up
to codimension two, and that the Igusa constants of the $1$-, $2$- and
$3$-regular base divisors are $2$, $3$ and $4$, as claimed. However, the proofs
of the main theorems (Theorems 6.1, 7.2 and 8.3) are not complete. They rest on
Hales' reduction lemma (Lemma 5.19), whose hypothesis cannot be met for the
cuspidal endoscopic groups of $G_2$. The hypotheses on the support of $f$ in
those theorems can never be satisfied. The endoscopic group $PGL(3)$ is given
the wrong roots. Several elliptic tori that carry nontrivial $\kappa$ are
missing. There are also many misprints, notational clashes and unproved
assertions. I recommend a major revision before the manuscript is circulated
as a proof of transfer for $G_2$.
\end{abstract}

\tableofcontents

\section{What I read, and how I checked it}

\paragraph{Files read.} I read all fifteen uploaded files completely:
\texttt{g2new0.tex} through \texttt{g2new9.tex}, the variants
\texttt{g2new5d.tex} and \texttt{g2new9d.tex}, \texttt{g2newapp.tex},
\texttt{g2newref.tex}, and the house style guide
\texttt{joyner-math-paper-style-guide.md}. I also read the text that follows
\verb|\enddocument| in \texttt{g2new2}, \texttt{g2new4}, \texttt{g2new6} and
\texttt{g2newapp}. That text was never typeset, but it bears on some issues
below. I refer to it as ``untypeset text''.

\paragraph{Sources consulted.} I could not consult Hales' thesis [H1], Hales'
$U(3)$ paper [H3], Langlands [L], Langlands--Shelstad [LS0]--[LS3],
Kempf et al.\ [K], Spaltenstein [Spal] or any other cited source in full.
I checked bibliographic data against online records (EuDML, Numdam, the IAS
archive of Langlands' papers, and a reference list in an AMS journal). I
consulted Hales' paper [H2] (Ast\'erisque 171--172) only through a
machine-generated summary of its Numdam PDF. Where I rely on that summary I say
so.

\paragraph{Verification levels.} Following \S7 of the style guide, every
point carries a level.
\Lv{1}: an independent derivation by me, by hand or by an exact computation
whose script is attached.
\Lv{2}: a check against a source I actually consulted.
\Lv{3}: reliance on memory of the literature.
Every issue also carries a verdict: \emph{correct}, \emph{misprint},
\emph{fixable gap} or \emph{false}.

\paragraph{Computations.} The following Python scripts accompany this report.
They use exact integer and symbolic arithmetic (SymPy), and they run in a few
minutes.
\begin{itemize}
\item \texttt{g2\_toric\_divisors.py}: the toric structure of $Z$, its facets,
all branches of $\lambda=0$ with $a(E)$, $b(E)$, $\beta(E)$, and which patches
meet which divisors.
\item \texttt{g2\_saturation\_full.py}: a complete proof that the semigroup of
$Z$ is saturated, so that $Z$ is normal.
\item \texttt{g2\_patch\_jacobians.py} and \texttt{g2\_patch\_regularity.py}:
the patch formulas of \S3, the Igusa forms of \S4, and the forms (8.1) and
(8.4).
\item \texttt{g2\_exceptional\_divisors.py}: the singular torus orbits of $Z$
and candidate exceptional divisors of a resolution.
\item \texttt{g2\_tori\_endoscopy.py}: $H^1(F,T)$ for every elliptic torus of
split $G_2$, and the root system of the endoscopic group for every $\kappa$.
\end{itemize}

\section{Summary of the manuscript}

Section 0 states the goal. The goal is to prove the matching of Shalika germs
for the $r$-regular unipotent orbits of $G_2$, $r=1,2,3$, by Langlands' Igusa
method. Together with Shelstad's result for $r=0$ and Rogawski's for the
trivial orbit, this is said to give the smooth matching (local transfer)
theorem for $G_2$. Section 1 recalls the variety of stars, Langlands'
coordinates $z(W,\rho)$, the partial resolution $S_1$ and the varieties $Y''$,
$Y_1$ and $Y_C$. Section 2 computes the Galois action on coordinates.
Section 3 introduces the monomial variety $Z$ and the ``fundamental'' and
``essential'' patches, and claims that $Z$ is normal and Cohen--Macaulay.
Section 4 lists the patches for $G_2$ and computes Igusa constants. Section 5
recalls Igusa's expansion, proves a ``no spurious divisors'' lemma and states
Hales' reduction lemma. Sections 6, 7 and 8 treat the $1$-, $2$- and
$3$-regular germs, with endoscopic groups $PGL(3)$ and
$(SL(2)\times SL(2))/\{\pm1\}$. Section 9 states a general conjecture. The
appendix computes a Chevalley basis via triality, the cuspidal endoscopic
groups, $H^1(\Gamma,T)$ for some tori, and the stability of the $2$- and
$3$-regular orbits.

The strategy has three steps. First, identify the unique base divisor that
contributes to each germ. Second, show that its Igusa coefficient depends on
the tangent direction $X$ only through a product of roots of $H$. Third, apply
Hales' reduction lemma to avoid computing that dependence exactly. The third
step does not work for $G_2$ (Issue~\ref{iss:reduction}), and the second is
asserted rather than computed (Issue~\ref{iss:pX}).

\section{Overall assessment}

The manuscript contains real and useful work. In particular, it gives an
explicit treatment of the Igusa variety of $G_2$, the essential patches, and
the identification of the base divisors $E_{\emptyset,\alpha,\delta}$ and
$E_{\emptyset,\beta,\zeta}$ that carry the $2$- and $3$-regular germs. My
independent computations confirm the core geometric numbers (Section~\ref{sec:correct}).

The main theorems are not established, for four independent reasons.
\begin{enumerate}
\item The descent step, Lemma 5.19, cannot be applied to the cuspidal
endoscopic groups of $G_2$ (Issue~\ref{iss:reduction}).
\item The hypotheses of Theorems 6.1, 7.2 and 8.3 are never satisfied
(Issue~\ref{iss:support}).
\item The roots of the endoscopic group $PGL(3)$ are misidentified
(Issue~\ref{iss:roots}).
\item Some elliptic endoscopic data are missing, so the smooth matching
theorem would not follow even if the three theorems were proved
(Issue~\ref{iss:tori}).
\end{enumerate}
The claim in \S0 that the paper ``proves the smooth matching theorem for
$G_2$'' should be withdrawn or made conditional.

For context in a revision (\Lv{3}): local transfer for $p$-adic groups is now
known in general. Waldspurger showed that the fundamental lemma implies
transfer (``Le lemme fondamental implique le transfert'', Compositio Math.\
105 (1997)), and Ng\^o proved the fundamental lemma (Publ.\ Math.\ IH\'ES 111
(2010)). A revision should therefore present the paper as an explicit
Igusa-theoretic computation of the $G_2$ germs, not as the only route to
transfer.

\section{Major issues}

\begin{issue}[Hales' reduction lemma does not apply]\label{iss:reduction}
\loc{\S0 (the paragraph on $\Delta_0(X)$), Lemma 5.19, and the ends of the
proofs of Theorems 6.1, 7.2 and 8.3}
Lemma 5.19 needs $Y$ in the Lie algebra of the center of the endoscopic group
$H$ with $\operatorname{Cent}(Y,G)=H$. Both cuspidal endoscopic groups of
$G_2$, namely $PGL(3)$ and $(SL(2)\times SL(2))/\{\pm1\}$, are semisimple of
rank two, the rank of $G_2$. So $\operatorname{Lie}(Z(H))=0$. Moreover, the
centralizer of a semisimple element $Y$ of $\mathfrak t$ is a Levi subgroup.
The Levi subgroups of $G_2$ are $T$, the two $GL(2)$'s and $G_2$. None of them
is $SL(3)$ or $SO(4)$, which are the subgroups of $G_2$ the paper uses as
``$H\subset G$''. So no admissible $Y$ exists. Read literally, the hypothesis
is vacuous, and the lemma would give transfer with no information about the
germs at all. That cannot be the intended statement.

The machine summary of [H2] that I consulted supports this reading. Hales'
quasi-split reduction (Prop.\ 2.11 of [H2]) descends to elements $s$ of the
connected center torus $S$ of $H$ close to the identity. For the cuspidal
$SO(4)$ case of $GSp(4)$, Hales instead computes the dependence on the tangent
direction explicitly: the germ is $|T_1-T_2|\,c(T,\kappa)$ (\S6 of [H2]).
The $G_2$ case needs the same kind of explicit computation, which \S0 says the
descent ``remarkably'' avoids.

\verdict{false as applied; the proofs of Theorems 6.1, 7.2 and 8.3 are
incomplete}
\level{\Lv{1} for the Lie-theoretic facts; \Lv{2} (machine summary only) for
[H2]; \Lv{3} for [H3, Lemma 6.1], which I could not consult}

\fix{Compute the $X$-dependence of the full coefficient
$\kappa(t_*)\,|\omega_{E(u)}|$, including the $\kappa$-factor, as in [H2,
\S6]. Then compare it with the $H$-side directly.}
\end{issue}

\begin{issue}[The support hypotheses are never satisfied]\label{iss:support}
\loc{Theorems 6.1, 7.2 and 8.3; the decomposition $f=f_1+f_2$ before Lemma
5.19; the conclusion of Lemma 5.19; Conjecture 9.1; the text of \S\S6--8}
Each theorem assumes that the support of $f$ meets the $r$-regular orbit
$\cO_r$ and no other unipotent orbit. The support of $f\in C_c^\infty(G(F))$
is open. Every unipotent $u\in G(F)$ lies in the unipotent radical $N$ of an
$F$-Borel subgroup (Borel--Tits, $G_2$ being split). Regular unipotents form
a Zariski-dense open subset of $N\cong\mathbb A^6$, so they are $p$-adically
dense in $N(F)$. Hence every neighborhood of $u$ meets the regular orbits. An
$f$ whose support meets $\cO_r$ also meets $\cO_0$ and every orbit between
them. The hypotheses are therefore never satisfied, and the theorems are
vacuous as stated. The same objection applies to $f_2$ in the decomposition
before Lemma 5.19. The phrase ``orbits outside the Zariski closure of $\cO$''
names the \emph{larger} orbits, which $f_2$ cannot avoid.

\verdict{false as stated; fixable}
\level{\Lv{1}}

\fix{State the theorems as germ identities. Alternatively, use test functions
with $\mu_{\cO'}(f)=0$ for every unipotent class $\cO'\neq\cO$. Such functions
exist because the unipotent orbital integrals are linearly independent.}
\end{issue}

\begin{issue}[The roots of the endoscopic group $PGL(3)$]\label{iss:roots}
\loc{\S8 (the definition of $\Sigma_{H_1}^+$, (8.4) and (8.7)); \S4, Lemma
4.4; Appendix A.5; untypeset text after \S6, (6.8)}
The paper takes $\Sigma_{H_1}^+=\{\beta,\epsilon,\zeta\}$, the long roots.
That is the root system of the \emph{pseudo-Levi subgroup} $SL(3)\subset G_2$.
The \emph{endoscopic group} $H_1=PGL(3)$ is different. Its roots, transported
to $T_G$, are the roots $a$ of $G$ with $\kappa(a^\vee)=1$. For the cubic
torus (Galois image $\langle(s_\beta s_\alpha)^2\rangle$) we have
$H^1(F,T)=X_*/(\sigma-1)X_*\cong\mathbb Z/3$. In the coroot basis
$(\alpha^\vee,\beta^\vee)$, $(\sigma-1)X_*$ is spanned by $(1,0)$ and $(0,3)$.
A nontrivial $\kappa$ is trivial exactly on the coroots
$\alpha^\vee=(1,0)$, $\gamma^\vee=(1,3)$ and $\delta^\vee=(2,3)$. So
\[
\Sigma_{H_1}^+=\{\alpha,\ \alpha+\beta,\ 2\alpha+\beta\},
\]
the \emph{short} roots. Equivalently, $\widehat H_1=SL(3)$ is generated by the
long roots of $\widehat G=G_2(\mathbb C)$, which are the coroots of the short
roots of $G$. The set $\Delta_2^\vee=\{\beta^\vee,(3\alpha+\beta)^\vee\}$ used
in A.5 consists of short roots of $\widehat G$, and they do not span a closed
subsystem: $\beta^\vee+\zeta^\vee=\gamma^\vee$ is a root.

For $H_2$ the paper's $\Sigma_{H_2}^+=\{\alpha,\zeta\}$ is correct, because
there the pseudo-Levi and the endoscopic root systems agree.

A short root system can still parametrize $\cO_3$ through the long-root
pseudo-Levi, so the coordinates in \S8 need not change. But (8.7), and the
claim after it that $p_{\cO_3,T}$ is a product over roots ``in $H$'', must use
the short roots, because $D_{H_1}$ and the transfer factor do.

\verdict{false; fixable}
\level{\Lv{1} (script \texttt{g2\_tori\_endoscopy.py})}
\end{issue}

\begin{issue}[Missing elliptic endoscopic data]\label{iss:tori}
\loc{Appendix A.6 (``All other non-split Cartans \dots do not contribute''),
Theorems 6.1, 7.2 and 8.3, the proof of Lemma 5.17, and \S0}
I computed $H^1(F,T)=\hat H^{-1}(\Gamma,X_*)$ for every elliptic Galois image
$\Gamma\subset W(G_2)$:

\smallskip
\begin{center}
\begin{tabular}{lll}
\toprule
Galois image in $W$ & $H^1(F,T)$ & roots of $H$ for $\kappa\neq1$\\
\midrule
$\{\pm1\}$ (quadratic $T$) & $(\mathbb Z/2)^2$ & $\{\alpha,\zeta\}$, $\{\beta,\delta\}$, $\{\gamma,\epsilon\}$\\
$\langle s_\alpha,s_\zeta\rangle$ and conjugates (biquadratic) & $\mathbb Z/2$ & one orthogonal pair\\
cyclic of order $3$ & $\mathbb Z/3$ & $\{\alpha,\gamma,\delta\}$\\
$S_3$ containing the short reflections & $\mathbb Z/3$ & $\{\alpha,\gamma,\delta\}$\\
cyclic of order $6$; $S_3$ with long reflections; $W$ & $0$ & --- \\
\bottomrule
\end{tabular}
\end{center}
\smallskip

The quadratic tori and the short-reflection $S_3$ tori carry nontrivial
$\kappa$ and are omitted. The proof of Lemma 5.17 also asserts that
$H^1(\Gamma,T_G)$ is cyclic of order $m\in\{2,3\}$, which fails for quadratic
tori.

The conclusions $H^1=\mathbb Z/2$, $0$ and $\mathbb Z/3$ in A.6 are right, but
the computations given for them are not.
\begin{itemize}
\item The displayed formula
$s_\beta s_\alpha(m\alpha^\vee+n\beta^\vee)=(2m-n)\alpha^\vee+(m-n)\beta^\vee$
has determinant $-1$, so it cannot be a rotation. The correct formula is
$(n-m)\alpha^\vee+(2n-3m)\beta^\vee$.
\item The ``cubic'' formula has determinant $3$.
\item The two coboundary identities are wrong.
\item The biquadratic computation takes a quotient of $H^1$ groups, which is
not valid. The correct computation is $(s_\alpha-1)X_*+(s_\zeta-1)X_*=\mathbb Z\alpha^\vee+\mathbb Z\zeta^\vee$, of index $2$.
\end{itemize}

\verdict{false (the claim of completeness); misprints in the computations}
\level{\Lv{1} (script)}
\end{issue}

\begin{issue}[Which divisors contribute]\label{iss:divisors}
\loc{\S4 (the table (4.3) and the dismissal of exceptional divisors), \S6
(``three base divisors with $\beta=2$''), Lemma 0.8, Lemma 5.14, \S\S7--8}
The toric computation lists every branch of $\lambda=0$ on $Z$:

\smallskip
\begin{center}
\begin{tabular}{lccccl}
\toprule
branch & $a$ & $b$ & $\beta$ & image & patches meeting it\\
\midrule
$E_{\{\alpha,\beta\}}$ & 1 & 1 & 1 & regular & $Z_{\{\alpha,\beta\}}$\\
$E_{\alpha}=E_{\emptyset,\alpha,\gamma}$ & 1 & 2 & 2 & subregular & $Z_\alpha$, $Z_{\emptyset,\alpha,\gamma}$\\
$E_{\beta}=E_{\emptyset,\beta,\gamma}=E_{\emptyset,\beta,\delta}=E_{\emptyset,\beta,\epsilon}$ & 1 & 2 & 2 & subregular & $Z_\beta$, $Z_{\emptyset,\beta,\gamma/\delta/\epsilon}$\\
$E_{\emptyset,\alpha,\delta}$ & 2 & 6 & 3 & $2$-regular & $Z_{\emptyset,\alpha,\delta}$\\
$E_{\emptyset,\alpha,\epsilon}=E_{\emptyset,\alpha,\zeta}$ & 3 & 11 & $11/3$ & --- & $Z_{\emptyset,\alpha,\epsilon}$, $Z_{\emptyset,\alpha,\zeta}$\\
$E_{\emptyset,\beta,\zeta}$ & 2 & 8 & 4 & $3$-regular & $Z_{\emptyset,\beta,\zeta}$\\
$E_{\emptyset}$ & 1 & 7 & 7 & trivial & $Z_\emptyset$\\
\bottomrule
\end{tabular}
\end{center}
\smallskip

\noindent Four points follow.
\begin{enumerate}[label=(\alph*)]
\item $E_{\emptyset,\alpha,\gamma}$ is the same divisor as $E_\alpha$. So
there are two base divisors with $\beta=2$, not three. The redundancy remark
in \S4 covers $\beta$ but not $\alpha$.
\item On $Z_{\emptyset,\alpha,\epsilon}$ the form is
$-x(\epsilon)^5z(\beta)z(\alpha)^{10}w(\epsilon)^{-6}$, not
$x(\epsilon)^5z(\beta)z(\alpha)^9w(\delta)^{-6}$. So $b=11$ and
$\beta=11/3$, not $10/3$, and the table entry $7/3$ should be $8/3$.
$E_{\emptyset,\alpha,\epsilon}$ and $E_{\emptyset,\alpha,\zeta}$ are one
divisor.
\item A divisor with $\beta=11/3$ would contribute $\theta(\lambda)|\lambda|^{8/3}$,
a term that cannot occur in the germ expansion. Its coefficient must therefore
vanish, and the paper should show this, for example by a ramification argument
for $\theta$ with $\theta^3=\kappa_E$. The manuscript never mentions this
divisor's contribution.
\item The dismissal of exceptional divisors (``because of lemma 0.8'') is
circular, since Lemma 0.8 concerns strict transforms only. $Z$ is singular
along five torus-orbit closures (listed in the output of
\texttt{g2\_exceptional\_divisors.py}). Any resolution adds rays $u$ in the
cone $\sigma$, and their Igusa constant is $\langle m_\omega,u\rangle/\langle\lambda,u\rangle$.
Integral values cannot be excluded in advance. For example,
$u=(3,2,2,2,3,5,7)$ lies in a singular face and has $\beta=2$. Whether the
$\beta=2,3,4$ coefficients receive extra contributions therefore depends on
the resolution chosen, and on the $F$-points and images of its exceptional
divisors. The spurious divisors over $Y_1\setminus Y''$ are likewise not
treated for $r=2,3$.
\end{enumerate}

The proof of Lemma 5.14 also has gaps. The homogeneity degrees ``by (1.22)''
are asserted, not computed. Only the stratum $D=E\cap E'$ is considered,
whereas Lemma 5.7 gives $F_1$ as a sum over strata. The contradiction requires
the $D$-term to be nonzero. And $e=1$ alone does not show that spurious
divisors fail to contribute.

\verdict{fixable gaps; the value $\beta=10/3$ is false}
\level{\Lv{1} (scripts)}
\end{issue}

\begin{issue}[The $X$-dependence and the measures on orbits and fibers]\label{iss:pX}
\loc{(6.5a), (6.5b), (7.3)--(7.5), (8.1)--(8.7)}
The decisive quantities $p_{\cO,T}(X)$ in (7.5) and (8.7) are asserted to be
products over roots of $H$ ``using (1.22)''. The exponents are never computed,
and neither is the $X$-dependence of $\kappa(t_*)$. In addition:
\begin{itemize}
\item (7.3) uses $dz(\alpha)$, but $z(\alpha)=0$ on $E_{\emptyset,\alpha,\delta}$.
It lists six differentials for the $8$-dimensional orbit $\cO_2$, and it
includes $dv(\alpha)$ although $\alpha\in\Sigma_{H_2}$. The $Z$-coordinates
available on $\cO_2$ are $w(\gamma)$, $w(\delta)$, $x(\delta)$ and
$x(\epsilon)$, and they should be paired with $v(\beta)$, $v(\gamma)$,
$v(\delta)$ and $v(\epsilon)$. ``$\Sigma_{\cO_3}$'' in (7.3) should read
$\Sigma_{\cO_2}$.
\item In (8.1) and (8.5), $w(\delta)$ should read $w(\zeta)$. With that
correction, (8.1) agrees with my computation:
$\omega_E=x(\zeta)^2w(\zeta)^{-3}z(\alpha)^{-2}$. The display after (8.2)
shows $\omega_{Z}$, mislabeled as $\omega_E$.
\item For (8.4) I find
$dx(\alpha)\,dx(\gamma)\,dx(\delta)/\lambda^2=2x(\zeta)w(\zeta)^{-1}z(\alpha)^{-2}\,dz(\alpha)\,dw(\gamma)\,dw(\delta)$,
not $z(\alpha)^{-1}$. More fundamentally, defining a measure on $\cO$ as a
$\lambda\to0$ limit of ratios is not justified. It is better to use the
invariant measure on $\cO$, which is unique up to a constant [R], and to
compute $\omega_{E(u)}$ from the fibration.
\item In \S6, $E_\alpha$ yields the factor $|\beta(X)|$. For $H_2$, $\beta$ is
not a root of $H$. The untypeset text after \S6 records an unresolved
discrepancy of this kind.
\item The translation invariance of $\kappa(t_*)$ ``by continuity'' needs a
bound uniform in $X$, which is not given.
\item Transfer needs one $f^H$ for all $(T,\kappa)$ attached to $H$. The
theorems fix $T_G$ before choosing $f^H$.
\end{itemize}

\verdict{fixable gaps (substantial)}
\level{\Lv{1}}
\end{issue}

\begin{issue}[Lemma 3.5, Lemma 4.4 and the ``regular branch'']\label{iss:richardson}
\loc{Lemma 3.5; \S3 (the regular branch $E_\emptyset$); Lemma 4.4}
The convention ``Hales' $E_S$ is my $E_{\Delta-S}$'' has not been carried
through. Three statements are affected.
\begin{itemize}
\item As printed, Lemma 3.5 gives for $S_1=\Delta$ the generators $\lambda$,
$x(\alpha)$, $x(\beta)$, $w(\gamma),\dots,w(\zeta)$. Then
$x(\gamma)=w(\gamma)x(\alpha)x(\beta)/\lambda$ is not regular. The correct
generators are $\lambda$, $z(\alpha)$ for $\alpha\in\Delta-S_1$, $x(\alpha)$
for $\alpha\in S_1$, $x(\rho)$ for $\rho\notin\Sigma^+_{\Delta-S_1}$, and
$w(\rho)$ for $\rho\in\Sigma^+_{\Delta-S_1}-\Delta$. The generator count then
equals $p+1$.
\item Lemma 4.4 should read $\beta(E_S)=1+|\Sigma^+_{\Delta-S}|$, with the
parabolic $P_{\Delta-S}$. As printed, it gives $\beta(E_\Delta)=7$, against
the value $1$ in the table.
\item ``The regular branch $E_\emptyset$ satisfies $\pi(E_\emptyset)=V$'' is
wrong in the paper's own convention. $E_\emptyset$ is the trivial-orbit branch
($\beta=7$), and the regular branch is $E_\Delta$.
\end{itemize}
The lists of generators in \S4 are correct. I checked all twelve patches.

\verdict{false as printed; misprint-level repair}
\level{\Lv{1}}
\end{issue}

\begin{issue}[The definition of $S''$ and $Y''$]\label{iss:Spp}
\loc{\S0, Step 2; Definition 1.16}
$S''$ is defined as ``the image of $S_0$ in $S_1$'' (and in Definition 1.16
as the image ``in $Y_1$''). A point of $S_0$ is a regular star, so every
$z(W,\rho)\neq0$ there, and over $\lambda=0$ the only element of
$\bigcap_W B(W)$ that is unipotent is $1$. The base divisors, on which
$z(\alpha)=0$, would then not lie in $Y''$. Presumably $S''$ is meant to be the
open subset of $S_1$ on which all $z_1(W,\rho)\neq0$, with $z(\alpha)$
arbitrary, as in Hales' thesis.

\verdict{fixable gap}
\level{\Lv{1} for the inconsistency; \Lv{3} for the intended definition}
\end{issue}

\begin{issue}[Lemma 3.7 and the singular locus]\label{iss:normal}
\loc{Lemma 3.7; \S3 (the claim that the singular locus lies in $K\cup K_0$);
Conjecture 3.4}
The proof of normality asserts saturation without checking it.
\begin{itemize}
\item For $G_2$ I proved saturation by an exhaustive exact computation, which
covers every lattice point of the fundamental parallelepipeds (degree at most
$6$). So $Z$ is normal for $G_2$. It is Cohen--Macaulay by Hochster's theorem
[Ho], which is in the bibliography but never cited. For general $G$ the claim
is unproved.
\item ``Set-theoretic complete intersection'' is not proved. The six equations
(3.1) also vanish on the $10$-dimensional set $\lambda=z(\alpha)=z(\beta)=0$,
so they do not cut out $Z$.
\item The singular-locus claim is correct, but the fifth patch is redundant.
On $Z$ one has $x(\epsilon)z(\alpha)^3=x(\beta)w(\epsilon)$, so that patch lies
inside $\{x(\beta)\ne0\}$.
\item $K$ contains the divisor $E_{\emptyset,\alpha,\delta}$, so ``beginning
with $K$ as the center of the first blow-up'' needs rewording.
\item ``The branches of $E_{Y''}$ are reduced'' contradicts the multiplicities
$2$ and $3$ found in \S4.
\end{itemize}
Conjecture 3.4 does hold for $G_2$: all essential patches are regular, and
every torus-invariant divisor of $Z$ meets one of them. The text proves only
the weaker statement about $K\cup K_0$.

\verdict{fixable gaps; the claims are true for $G_2$}
\level{\Lv{1} (scripts)}
\end{issue}

\begin{issue}[Section 1 foundations]\label{iss:sec1}
The following points concern \S1.
\begin{itemize}
\item \emph{The map $\phi$ is not well-defined.} A regular conjugacy class
meets $T_G$ in a full Weyl orbit, so ``the point in $T_G$ which meets the
conjugacy class of $g$'' is defined only in $T_G/\Omega$. The star must be used
to choose the point. (\Lv{1}; fixable.)
\item \emph{Negation error.} Not lying in $S'(B_+,B_-)$ means that for
\emph{some} $\alpha\in\Delta$, $z=0$ for \emph{all} $\omega$. The text says
``for all $\alpha$ and $\omega$'' and concludes that the star lies on the
diagonal. The intended conclusion, that such stars are not regular, still
follows from Lemma 1.12. (\Lv{1}; fixable.)
\item $\rho=\omega\alpha$ (in the definition of a pair and of $S'$)
contradicts $\alpha_{i+1}=\omega_i\rho_{i+1}$; it should read
$\rho=\omega^{-1}\alpha$. The sign in (1.8) is opposite to that in (1.9) and
(1.10). In the definition of a root, $t^{-1}\theta_\rho(\xi)t$ should be
$t\theta_\rho(\xi)t^{-1}$. ``$\ker\rho\subset X^*(T)\otimes\mathbb R$'' should
be the hyperplane $\rho^\perp$. (\Lv{1}; misprints.)
\item \emph{Equation (1.22).}
\begin{itemize}
\item The $c_1$-term of $w(\gamma)$ lacks the factor
$z(\sigma_\beta^{-1}W_0,\alpha)$ that appears in $x(\gamma)$.
\item The last term of $w(\zeta)$ lacks $c_7$ and has $n_\zeta$ for $n_\delta$.
\item The constants $c_1,\dots,c_7$ and the meaning of ``higher order terms in
$\lambda$'' are never specified.
\item The final list gives seven ``algebraically independent'' coordinates.
Together with $\lambda$ that makes eight functions on the $7$-dimensional
patch $Y_1(B_+,B_-)$, since $z(W_0,\alpha)=1/n_\alpha$ and
$z(W_0,\beta)=1/n_\beta$ already account for the star. The list needs
correcting.
\end{itemize}
(\Lv{1}; fixable.)
\end{itemize}
\end{issue}

\begin{issue}[Section 2 computations]\label{iss:sec2}
The following points concern \S2.
\begin{itemize}
\item The requirement $g^{-1}\sigma(g)\in T_G(\ovF)$ for all $\sigma$ would
make $B_G$ defined over $F$, which is impossible for elliptic $T_G$. It should
be $g^{-1}\sigma(g)\in N(T_G)(\ovF)$, and this gives $\sigma_T$. ``Factors
through $Gal(\ovF/E)$'' should read ``through $Gal(E/F)$''. (\Lv{1}.)
\item In (2.8) and (2.9), $b'$ should be
$c_\sigma\sigma(n_{\sigma_T^{-1}W_0})\,\sigma(b)\,\sigma(n_{\sigma_T^{-1}W_0})^{-1}c_\sigma^{-1}$.
The printed formula has $b$ for $\sigma(b)$, $c_\sigma$ in place of
$c_\sigma^{-1}$, and a double inverse. (\Lv{1}; misprint.)
\item In $SL(2)$, $\theta_\alpha(x)^{\theta_{-\alpha}(z)}=\theta_{-\alpha}\!\big(\tfrac{-xz^2}{1+xz}\big)\,\alpha^\vee(1+xz)\,\theta_\alpha\!\big(\tfrac{x}{1+xz}\big)$.
So $\phi^\alpha_\alpha=x(\alpha)/(1+x(\alpha)z)$, not $x(\alpha)$. The
analogous $y_\beta$ is $-xz^2/(1+xz)$. The torus factor also rescales the
other coordinates. (\Lv{1}.)
\item In $\theta_\beta(r)\theta_\alpha(s)$ the exponents are swapped. The
correct terms are $\theta_\delta(rs^2)\theta_\epsilon(rs^3)\theta_\zeta(r^2s^3)$,
because $i\alpha+j\beta$ carries $s^ir^j$. In $\phi^\alpha_\delta$, $c_2$
should be $c_1$. In the sentence after Lemma 2.10, $\theta_\rho(z)$ should be
$\theta_{-\rho}(z)$. In (2.17) the subscripts $\eta_{\alpha,\cdot}$ should be
$\eta_{\beta,\cdot}$. (\Lv{1}; misprints.)
\item The figure of the Dynkin curve is missing: the source has empty displays.
\end{itemize}
\end{issue}

\begin{issue}[The appendix]\label{iss:app}
The following points concern the appendix.
\begin{itemize}
\item \emph{A.2.}
\begin{itemize}
\item $k\times j=i$ should be $k\times j=-i$.
\item The form $\tfrac12(xy^*+y^*x)$ is not scalar-valued: for quaternions with
$x=1$ and $y=i$ it equals $-i$. The correct form is $\tfrac12(xy^*+yx^*)$.
\item The triality theorem holds for proper orthogonal maps (rotations), not
for arbitrary ones.
\item The ``proof'' checks only the infinitesimal identity, which does not
imply the global statement. A reference such as Springer--Veldkamp would
serve.
\end{itemize}
(\Lv{1}; \Lv{3} for the reference.)
\item \emph{A.4.} The matrix of $X_{\alpha+\beta}$ has entry $-1$ at position
$(1,3)$, while the note says $E_{13}+E_{24}+E_{25}$. All other matrices check
out (see Section~\ref{sec:correct}). (\Lv{1}.)
\item \emph{A.5.}
\begin{itemize}
\item The displayed basis of $P$ is wrong: $(1/2,0)\cdot\beta^\vee=-1/2$. The
correct basis is $(1/2,\sqrt3/2)$, $(0,\sqrt3)$. The conclusion $P=Q$ is right.
\item The $A_1\times A_1$ subsystem should be generated by $\alpha^\vee$ and
$\zeta^\vee$, not $\beta^\vee$ and $\zeta^\vee$, which are not orthogonal.
\item ``$Q^\vee$ strictly contains $P^\vee$'' should read ``is strictly
contained in $P_1^\vee$''.
\item The argument identifying $H_2$ is hard to follow, though the answer is
right: the character lattice $Q$ is the index-$2$ sublattice of $P_1$
containing $(\alpha+\zeta)/2$.
\item For $H_1$, see Issue~\ref{iss:roots}.
\end{itemize}
(\Lv{1}.)
\item \emph{A.7.} The stability arguments for $\cO_2$ and $\cO_3$ are
plausible. I confirmed $\mathfrak g_2(2)=F X_{\alpha+\beta}$ and
$\mathfrak g_2(2)=F X_\beta$ and the subregular grading. Two points need
attention.
\begin{itemize}
\item For $\cO_2$ the grading is odd, so the use of Rao's lemma should be
justified.
\item The claim ``$7$ $F$-classes'' in the stable subregular class conflicts
with my count. The component group is $A(u)=S_3$ and $H^1(F,G_2)=1$, so the
classes correspond to étale cubic $F$-algebras. There are $8$ of these when
$q\equiv1\pmod 3$ and $6$ when $q\equiv2\pmod 3$ ($p\neq2,3$). Please check
Assem's statement.
\end{itemize}
(\Lv{1} with \Lv{3} inputs.)
\item \emph{Untypeset weighted Dynkin diagrams.} The labels for the $2$- and
$3$-regular classes are swapped. The long-root class $A_1$ ($3$-regular) has
weights $(\alpha,\beta)=(0,1)$, and the short-root class $\tilde A_1$
($2$-regular) has $(1,0)$. (\Lv{1}.)
\end{itemize}
\end{issue}

\section{Minor issues, misprints and inconsistencies}

All entries are \Lv{1} unless marked otherwise.

\begin{longtable}{p{0.16\textwidth}p{0.56\textwidth}p{0.18\textwidth}}
\toprule
Location & Issue & Verdict\\
\midrule
\endhead
\S0 & ``$T'_G\cap\{1\}$'' should be $T'_G\cup\{1\}$. & misprint\\
\S0 & Property (e) sets $E=\pi^{-1}(u)$, a fiber. It should concern a divisor $E$ with $\pi(E)=\overline\cO$. & misprint\\
\S0 & The conjecture mixes $\Delta(\gamma_H,\gamma_\lambda)$ with $\Phi(\gamma_G,f)$. & misprint\\
\S0 & ``apply lemma 5.18'' should read ``lemma 5.19''. (5.18) is an equation. & misprint\\
\S0 & $D_H$ is used without absolute values in (0.1) and in Formula 0.6. & misprint\\
\S0 & ``Langland's'' should be ``Langlands's''. & spelling\\
\S0, \S3 & ``Cohen-Macauley'' should be ``Cohen--Macaulay'', here and in [Ho]. & spelling\\
\S0 & Lemma 0.8 says it is ``proven in \S\S4--5''. The equality of the principal values over $E$ and $\tilde E$ is never proved there. & fixable gap\\
\S0 & ``desingularize $Y''$ by blowing up four times'': the derivation is omitted, so the claim cannot be checked. & unverified\\
\S1 & ``Lemma 1.2 below'' does not exist. Tag (1.3) is missing. & misprint\\
\S1 & [Sp3] is cited but missing from the bibliography. & misprint\\
\S1 & Lemma 1.21: the formula for $c_{\rho_1,\dots,\rho_m}$ is not a precise statement. & readability\\
\S2 & ``the projection map $\pi$ of \S1.1'': there is no \S1.1. $X^0$ should be $X_0$. & misprint\\
\S3 & In g2new3 the macro \verb|\off| means $\OF$, so the first line prints ``over $\OF$''. $\ovF$ was probably intended. & misprint\\
\S3 & $x(\beta)\neq0$ patch: $w(\zeta)=x(\alpha)x(\zeta)z(\alpha)^4/x(\beta)^2$. & misprint\\
\S3 & $x(\gamma)\neq0$ patch: the exponent of $x(\gamma)$ in $w(\zeta)$ is $4$. & misprint\\
\S3 & $z(\alpha)\neq0$ patch: $w(\zeta)=x(\zeta)w(\epsilon)z(\beta)/x(\epsilon)$, with no $z(\alpha)^{-3}$. & misprint\\
\S3 & Lemma 3.6: the statement has $\rho>\sigma$, the proof uses $\rho\ge\sigma$. In general $G$, roots incomparable with $\sigma$ are not covered. & misprint / gap\\
\S3 & ``known \dots for $C_n$ [H1], [H2]'': [H2] treats only $GSp(4)$. ``False for $D_n$'' is unsupported. & \Lv{2}/\Lv{3}\\
\S3 & [Sp] is not a bibliography label. & misprint\\
\S4 & In $\omega_{Y''(B_+)}$ there is a doubled $\wedge$. & misprint\\
\S4 & $\sigma=\epsilon,\zeta$ forms: the denominator should be $w(\epsilon)$, $w(\zeta)$, not $w(\delta)$. & misprint\\
\S4 & Lemma 4.4: the second case of $\Ind$ should read ``$3$-regular''. & misprint\\
\S4 & Observation (II): for strict transforms, $\beta$ is unchanged whether or not $E$ contains a center, because $Z$ is normal. & readability\\
\S5 & Lemma 5.10(c): the chain $\overline\cO=\cO\cup\cO_m$ is garbled. & misprint\\
\S5 & In the proof of Lemma 5.7, $\sum_{1\le i\le p}c_j/(j-1)!$ has the wrong index and a spurious $(j-1)!$. & misprint\\
\S5 & In the proof of Lemma 5.17, ``(5.10)'' refers to a lemma, and ``$\xi$ satisfies (5.2)'' is unexplained. & fixable gap\\
\S5 & Lemma 5.19: with $\Gamma^{T,\kappa}$ normalized by $\Delta$, the homogeneity exponent is $2(r-|\Sigma_H^+|)$, not $2r$. & misprint\\
\S6 & $E_{\emptyset,\alpha,\gamma}=E_\alpha$; see Issue~\ref{iss:divisors}. & false\\
\S6 & (6.5b) drops the prefactor $|\lambda|\kappa(\cdot)D_H^{-1}$. & misprint\\
\S6 & ``$E$ is $E_\alpha$ or $E_\beta$, depending on $T_G$'' is unexplained, since both are $\beta=2$ divisors. & fixable gap\\
\S6 & In g2new6 \verb|\off| means $\ovF$, so the text prints $\exp(\ovF X)$. $\OF$ may have been intended. & misprint\\
\S7 & ``$\Sigma_{\cO_3}$'' in (7.3) should read $\Sigma_{\cO_2}$. The exponent $3/2$ of $|\lambda|$ is unexplained. & misprint\\
\S7 & $F$-rationality of $E_{\emptyset,\alpha,\delta}$ via observation (I) needs the uniqueness argument spelled out. & readability\\
\S8 & (8.6) should have $|\lambda|^3$, not $|\lambda|^2$; compare the display before (8.2). & misprint\\
\S8 & ``conjecture 3.4 (which we proved in \S3)'': \S3 does not prove it. & false (but true)\\
\S9 & Variant g2new9d: \verb|\off| means $\OF$, so (5) and (6) print $G(\OF)$, $H(\OF)$. [A] and [H4] are not in the bibliography. & misprint\\
App.\ A.6 & ``(a) simple reflections $\langle s_\alpha,s_\zeta\rangle$'': $s_\zeta$ is not simple. & misprint\\
Untypeset & The ``Lemma 9.1'' in the text after \S6 belongs to no section. ``Lemma 1.21'' in the text after \S2 duplicates Lemma 6.4. & housekeeping\\
Throughout & ``homogeniety'', ``throughtout'', ``neighborhoood'', ``occuring'', ``patched'' (for patches), ``Automomorphic'', ``Represenattions'', ``Birkhauser''. & spelling\\
\bottomrule
\end{longtable}

\section{Results I verified as correct}\label{sec:correct}

\begin{itemize}
\item \emph{Normality.} $Z$ is normal for $G_2$. The essential patches listed
in \S4 are all regular and cover $Z$ up to codimension two, so Conjecture 3.4
holds for $G_2$. (\Lv{1}.)
\item \emph{Igusa constants.} $\beta=1$ ($E_{\{\alpha,\beta\}}$), $2$ ($E_\alpha$,
$E_\beta$), $3$ with $a=2$, $b=6$ ($E_{\emptyset,\alpha,\delta}$), $4$ with
$a=2$, $b=8$ ($E_{\emptyset,\beta,\zeta}$), and $7$ (trivial orbit). The forms
on $Z_\alpha$, $Z_\beta$, $Z_{\emptyset,\alpha,\gamma}$,
$Z_{\emptyset,\alpha,\delta}$ and $Z_{\emptyset,\alpha,\zeta}$ agree with mine
up to the misprinted denominators. (\Lv{1}.)
\item \emph{Patch formulas.} The patch formulas of \S3 are correct apart from
the three misprints listed. (\Lv{1}.)
\item \emph{Orbit bookkeeping.} The orbit dimensions $12-2r$, the homogeneity
of the normalized germs ($|\lambda|^{2r}$), the relation $\beta-1=r$, and the
dimension counts $\dim\cO_2=8=14-\dim SO(4)$ and $\dim\cO_3=6=14-\dim SL(3)$
are correct. (\Lv{1}.)
\item \emph{Lemma 5.6} is correct. (\Lv{1}.)
\item \emph{Triality.} The root bookkeeping of A.3 and (A.11) is correct: the
fixed and permuted roots of $D_4$ and the formulas for $\phi(\epsilon_i)$. So
are all Chevalley matrices $X_\sigma$, $h_\sigma$ in A.4 apart from the sign
noted, and the torus $T'$ with $\alpha=a_3$, $\beta=a_2a_3^{-1}$. (\Lv{1}.)
\item \emph{$H_2$.} The endoscopic group $H_2=(SL(2)\times SL(2))/\{\pm1\}$
with $\Sigma_{H_2}^+=\{\alpha,\zeta\}$ is correct, and $H_1\cong PGL(3)$ as an
abstract group is correct. (\Lv{1}.)
\item \emph{Rational orbits.} The weighted gradings in A.7 are right, as is
the conclusion that $\cO_2$ and $\cO_3$ each consist of a single rational
orbit. (\Lv{1} for the gradings.)
\end{itemize}

\section{Readability}

\begin{enumerate}
\item \emph{Notation clashes.}
\begin{itemize}
\item $\beta$ is both the long simple root and the Igusa exponent $b/a$. For
example, ``$\beta(E_{\emptyset,\beta,\zeta})=4$''.
\item $\gamma$ is a root, a group element and a constant in (5.2).
\item $\delta$ is a root, a cohomology class, and the blow-up indicator in
(II).
\item $\epsilon$ is a root and the Mellin cutoff. $p$ is the prime, $|\Sigma^+|$
and an index in Lemma 5.7. $\pi$ is the projection and the uniformizer.
\item $\sigma$ is a Galois element, a root and the map $S_2\to\Sigma^+$.
$\alpha$ is also the unit in (5.1).
\item $E$ is a divisor and the splitting field. $S_1$ is a subset of $\Delta$
and the variety $S_1$. $K$ is a subvariety and a maximal compact subgroup.
\end{itemize}
I suggest renaming the Igusa exponent (for example $\nu(E)$) and the roots
$\gamma,\dots,\zeta$ (for example $\rho_3,\dots,\rho_6$).
\item The macro \verb|\off| means $\ovF$ in some files and $\OF$ in others.
This produces at least four printed errors.
\item Conventions are never fixed in one place: $B^g=g^{-1}Bg$, the meaning
of $\sigma_T$, which Borel contains the positive roots, the sign of $z$, and
the Chevalley normalization. The reader has to infer them from formulas, some
of which conflict (Issue~\ref{iss:sec1}).
\item The constants $c_i$, $d_i$, $e$, $f_i$ are ``computable'' but never
computed. The phrase ``higher order terms in $\lambda$'' is used for
expressions whose other coordinates also depend on $\lambda$.
\item Several key steps are ``left to the reader'', ``omitted'' or
``similar'': Lemma 5.6, the case $e=3$ in Lemma 5.14, the four-step blow-up,
and the details of (1.22).
\item The Dynkin curve figure is missing.
\end{enumerate}

\section{Bibliography check}

\begin{longtable}{p{0.12\textwidth}p{0.62\textwidth}p{0.16\textwidth}}
\toprule
Label & Finding & Level\\
\midrule
\endhead
{[Ro]} & Published as J. Rogawski, ``An application of the building to orbital integrals'', Compositio Math.\ 42 (1980) 417--423. The thesis title in the manuscript differs. & \Lv{2} (EuDML, Numdam)\\
{[S2]} & The title is ``A formula for regular unipotent germs'', Ast\'erisque 171--172 (1989) 275--277. & \Lv{2} (Numdam)\\
{[H2]} & The pages are 195--256. & \Lv{2} (Numdam)\\
{[H3]} & The pages are 303--333 (Univ.\ Montr\'eal, 1992). & \Lv{2} (AMS reference list)\\
{[L]} & The pages are 465--506. & \Lv{2} (IAS archive)\\
{[LS0]} & LNM 1041, pp.\ 250--279. The volume is dated 1984 (copyright 1983). & \Lv{2}\\
{[LS3]} & Correct as listed. & \Lv{2}\\
{[K]} & LNM 339 is dated 1973. The authors are Kempf, Knudsen, Mumford and Saint-Donat. & \Lv{3}\\
{[R]} & The pages should be checked: I recall 505--510. & \Lv{3}\\
{[SpSt]} & LNM 131 is from 1970. ``1986 (2nd ed.)'' should be checked. & \Lv{3}\\
{[Har]} & The year (1977) is missing. & \Lv{3}\\
{[Rep]} & This is part II (subregular term). Shelstad [S2] cites part I, PJM 113 (1984) 165--172. & \Lv{2}/\Lv{3}\\
missing & [Sp3], [Sp], [A], [H4], [D] are cited but have no entry. & \Lv{1}\\
uncited & [Ho], [Rep$'$], [Sp1] are never cited. [Ho] is the natural reference for Lemma 3.7. & \Lv{1}\\
\bottomrule
\end{longtable}

\section{Compliance with the house style guide}

The manuscript is a 1994 draft, so these items describe what a revision under
\S\S1--6 of the style guide will need. They are not criticisms of the
original.
\begin{itemize}
\item \emph{\S2 (structure).} The paper has no abstract, AI-disclosure
sentence, literature review, background section with examples, notation
table, summary of results or glossary.
\item \emph{\S1 (typesetting).} The source is AMS-TeX split over separate
files with hand-typed numbers and hand-typed cross-references. A revision needs
one amsart file, \verb|\label|/\verb|\ref|, theorem numbering within sections,
and citation keys of the form \texttt{Ha89}. The transcription delivered with
this report takes the first steps (single file, house-style keys, 12pt Times,
letter paper, a clean pdflatex build) but keeps the hand-typed numbering on
purpose.
\item \emph{\S4 (status of results).} Theorems 6.1, 7.2 and 8.3 should be
labeled conditional or conjectural until Issues~\ref{iss:reduction}--\ref{iss:pX}
are resolved. The same applies to Lemma 3.7 for general $G$, to Lemma 0.8, and
to Lemma 5.14.
\item \emph{\S5 (authorship).} The author line should give
\texttt{wdjoyner@gmail.com} in place of the USNA address. The acknowledgments
need the AI-usage statement.
\item \emph{\S6 (revisions).} The title, the subtitle ``Draft $n$ --- Corrected
and expanded version'' and the dates need to be added. A revision log is
needed, with a list of ``Remark 2026.$n$'' items for unverified claims.
\end{itemize}

\section{Notes on the transcription}

The single-file transcription follows \S6.7 of the style guide: it is a
literal transcription with no corrections. All errors listed above are kept.
Additions by me appear only as comments beginning ``Claude:''. Specific
choices:
\begin{enumerate}
\item The body follows \texttt{g2new5.tex} and \texttt{g2new9.tex}. Every
difference in the variants \texttt{g2new5d.tex} and \texttt{g2new9d.tex} is
transcribed in a comment or an \verb|\iffalse| block marked ``VARIANT''. I
could not tell which version is later. The ``d'' files align \S5 with the
wording of \S8 and extend Conjecture 9.1. They also cite [A] and [H4], which
are missing from the bibliography.
\item The text after \verb|\enddocument| is transcribed in \verb|\iffalse|
blocks marked ``UNTYPESET''. All of these blocks compile when switched to
\verb|\iftrue|.
\item The macro \verb|\off| is rendered as each file printed it, and doubtful
cases are flagged.
\item Missing citation labels are typeset literally.
\item Underlined bibliography titles are set in italics.
\item Long inline formulas were allowed to break after commas. This avoids
overfull boxes and changes no content.
\end{enumerate}

\section{Recommendation}

Major revision. The geometric core is worth preserving, and with the
corrections above it is largely right. A revised version should do five
things.
\begin{enumerate}
\item Replace the appeal to Lemma 5.19 by an explicit computation of the
$X$-dependence, including the $\kappa$-factor, in the style of [H2, \S6].
\item Restate the theorems as germ identities, or with orbital-integral
conditions on $f$.
\item Use the correct root system for $PGL(3)$, and treat the quadratic and
$S_3$ tori.
\item Account for the divisor with $\beta=11/3$ and for the exceptional and
spurious divisors.
\item Fix the misprints and the notation clashes.
\end{enumerate}
Until then, the main results should carry the label ``conditional'' (style
guide \S4.1).

\section*{Acknowledgments and AI disclosure}

This report was written by Claude (Opus 5.5, Anthropic) acting as referee at
the request of the author. The computations were run by Claude in Python with
SymPy and SciPy, and the scripts accompany the report. AI review is not an
independent certification of the results, and statements at level \Lv{3}
should be checked against the sources. Further details on the AI usage belong
in the revision log.

\end{document}
